\chapter{Pythonによる簡単なプログラミング}
\label{chap:python}
\chapterauthor{\AuthorOne}

\section{小さな関数を作る}
Pythonで数列の合計を計算してみましょう。
コード\ref{lst:python-total}は、渡された値を順番に加算します。
Pythonの実行環境は組版には不要です\footnote{コードは紙面に掲載するだけで、ビルド時には実行しません。}。

\begin{lstlisting}[caption={数列の合計を求める関数},label={lst:python-total}]
def total(values):
    result = 0
    for value in values:
        result += value
    return result

print(total([1, 2, 3, 4, 5]))  # 15
print("This deliberately long sample line demonstrates automatic wrapping within the code frame.")
\end{lstlisting}

\subsection{確認したいこと}
\begin{itemize}
  \item 関数の役割を短い文章で説明する。
  \item 入力と出力の例を用意する。
  \item 空のリストのような境界条件を確認する。
\end{itemize}

\section{実行結果を整理する}
表\ref{tab:python-results}に確認例を示します。
\begin{table}[htbp]
  \centering
  \caption{入力と期待する合計}
  \label{tab:python-results}
  \begin{tabular}{lr}
    \toprule
    入力 & 合計 \\
    \midrule
    \texttt{[1, 2, 3, 4, 5]} & 15 \\
    \texttt{[10, 20]} & 30 \\
    \texttt{[]} & 0 \\
    \bottomrule
  \end{tabular}
\end{table}

次の順で読み直すと、説明の不足を見つけやすくなります。
\begin{enumerate}
  \item コードの入力を確認する。
  \item 表とコードの出力を照合する。
  \item 第\ref{chap:algorithm}章の数式と対応付ける。
\end{enumerate}
Pythonの公式資料は\url{https://docs.python.org/ja/3/}を参照してください。
